\chapter{Inertial Sensors: Accelerometers, Gyroscopes, Magnetometers}

For thirty-three chapters we have treated signals as abstractions: a function $x(t)$,
a sequence $x[n]$, a vector in some space, a spectrum. Now we plug into the physical
world. The wrist tracker we are building toward does not receive a tidy waveform from
the heavens -- it has to \emph{generate} its signals by pressing tiny pieces of
silicon against the laws of mechanics and electromagnetism. This chapter is about the
three sensors that almost every motion-sensing project leans on: the
\textbf{accelerometer}, the \textbf{gyroscope}, and the \textbf{magnetometer}.
Together they form what the industry calls an \emph{IMU} (Inertial Measurement Unit),
and when a magnetometer is added, a \emph{MARG} array (Magnetic, Angular Rate, and
Gravity).

The goal here is not to memorize datasheets. It is to \emph{see}, geometrically, what
each sensor actually measures, what its raw signal looks like when a real arm moves,
and -- crucially -- what it cannot tell you no matter how cleverly you process it. By
the end you should be able to wire an MPU-6050 or BNO055 to a microcontroller, stream
its numbers, recognize the shape of a wrist flick in the trace, and understand exactly
why a wrist IMU sees the bulk motion of the hand but leaves the fingers to a different
sensor -- one we add in the capstone, not because fingers are unreachable, but because
they live in a different physical channel.

\section{Three sensors, three different physical questions}
\label{ch34:three-questions}

It is tempting to lump the three sensors together as ``the motion chip,'' because on
a modern board they share one package and one I2C address. But they answer three
genuinely different physical questions, and confusing them is the most common source
of bugs in motion code.

\begin{itemize}
\item The \textbf{accelerometer} answers: \emph{which way is the proof mass being
  pushed, and how hard?} It measures \emph{specific force} -- acceleration minus the
  effect of gravity, in a sense we will make precise -- in units of $\mathrm{m/s^2}$
  or $g$ (where $1\,g \approx 9.81\,\mathrm{m/s^2}$).
\item The \textbf{gyroscope} answers: \emph{how fast is the chip rotating right now?}
  It measures \emph{angular rate} about each of its three axes, in $\mathrm{deg/s}$ or
  $\mathrm{rad/s}$.
\item The \textbf{magnetometer} answers: \emph{which way is the ambient magnetic field
  pointing, as seen from the chip?} It measures the magnetic flux density vector in
  microtesla ($\mathrm{\mu T}$) or gauss ($1\,\mathrm{G}=100\,\mathrm{\mu T}$).
\end{itemize}

\begin{intuitionbox}
A clean mnemonic for what each sensor ``knows'':
\begin{itemize}
\item Accelerometer: knows \emph{down} (when you hold still) and \emph{shoves} (when
  you move). It is a tilt sensor that is also a bump detector.
\item Gyroscope: knows \emph{turning}. It is blind to which way is down and blind to
  steady motion -- it only feels change in orientation.
\item Magnetometer: knows \emph{north} (roughly). It is a compass, and like a compass
  it is easily fooled by nearby iron and motors.
\end{itemize}
Each sensor is excellent at one thing and useless at the others. Fusion
(Chapter~36) is the art of combining their strengths; this chapter is about
respecting what each one does and does not measure.
\end{intuitionbox}

\begin{figure}[H]
\centering
\begin{tikzpicture}[font=\small]
  \foreach \x/\name/\q/\unit in {0/Accelerometer/specific force/m$\cdot$s$^{-2}$,
       4.6/Gyroscope/angular rate/deg$\cdot$s$^{-1}$,
       9.2/Magnetometer/magnetic field/$\mu$T} {
     \draw[rounded corners,thick,fill=blue!4] (\x,0) rectangle (\x+3.8,1.5);
     \node[font=\bfseries] at (\x+1.9,1.15) {\name};
     \node at (\x+1.9,0.65) {\q};
     \node[font=\itshape\footnotesize] at (\x+1.9,0.2) {\unit};
  }
  \draw[rounded corners,thick,fill=violet!8] (3,-2.0) rectangle (9.8,-0.9);
  \node[font=\bfseries] at (6.4,-1.2) {IMU / MARG package (one I2C device)};
  \node[font=\footnotesize] at (6.4,-1.7) {3+3 axes, sometimes 3 more for the magnetometer};
  \foreach \x in {1.9,6.5,11.1} { \draw[-{Latex[]},gray] (\x,0) -- (6.4,-0.9); }
\end{tikzpicture}
\caption{The three inertial/field sensors and what each physically measures. On a
9-axis board (e.g.\ MPU-9250, BNO055) all three live in one package and share a
bus, but they answer independent physical questions.}
\label{ch34:fig-overview}
\end{figure}

\section{The accelerometer: a mass on springs}
\label{ch34:accel}

\subsection{The mechanical picture}

Strip away the silicon and an accelerometer is embarrassingly simple: a small
\emph{proof mass} suspended by springs inside a frame. When the frame accelerates, the
mass lags behind (inertia), the springs stretch, and the displacement of the mass
relative to the frame tells you the force. Newton's second law for the mass, written
in the frame of the chip, gives a spring displacement $x$ obeying
\[
  m\,\ddot{x} + b\,\dot{x} + k\,x \;=\; -\,m\,a_{\text{frame}} + m\,g_{\text{axis}},
\]
where $a_{\text{frame}}$ is the frame's acceleration and $g_{\text{axis}}$ is the
component of gravity along the sensing axis. In steady state ($\ddot{x}=\dot{x}=0$) the
displacement is proportional to $g_{\text{axis}} - a_{\text{frame}}$. That combination
is the whole story, and it has a name.

\begin{definition}[Specific force]
\label{ch34:def-specific-force}
An accelerometer measures \emph{specific force} (also called proper acceleration):
\[
  \vec{a}_{\text{meas}} \;=\; \vec{a}_{\text{true}} \;-\; \vec{g},
\]
where $\vec{a}_{\text{true}}$ is the true acceleration of the device in an inertial
frame and $\vec{g}$ is the gravitational acceleration vector (pointing \emph{down},
magnitude $\approx 9.81\,\mathrm{m/s^2}$). All quantities are expressed in the chip's
own (body) coordinate frame.
\end{definition}

This minus sign is the single most misunderstood fact about accelerometers, so let us
nail it with the two cases everyone gets wrong.

\begin{trapbox}
\textbf{A sitting accelerometer reads $+1\,g$ upward, not zero.}
A chip resting flat on a table is not accelerating ($\vec{a}_{\text{true}}=0$), so it
reports $\vec{a}_{\text{meas}} = -\vec{g}$, a vector of magnitude $9.81\,\mathrm{m/s^2}$
pointing \emph{up}. The table pushes the proof mass up against its springs exactly as
if the whole thing were accelerating upward at $1\,g$ in deep space. \\[4pt]
\textbf{A freely falling accelerometer reads zero.}
Drop the chip and $\vec{a}_{\text{true}} = \vec{g}$, so
$\vec{a}_{\text{meas}} = \vec{g}-\vec{g} = 0$. The proof mass and frame fall together;
the springs relax. This is why an accelerometer detects free fall and why phones can
sense a drop before impact.
\end{trapbox}

\begin{whybox}
Why does it measure $\vec a_{\text{true}}-\vec g$ rather than plain acceleration?
Because gravity acts on the proof mass and the frame \emph{equally} (equivalence
principle). Only forces transmitted \emph{through the springs} stretch them, and
gravity is not transmitted through the springs -- it pulls the mass directly. So the
springs report everything \emph{except} the gravity-driven part of the motion, which
is exactly $\vec a_{\text{true}} - \vec g$. The sensor cannot, even in principle,
distinguish ``sitting in gravity'' from ``accelerating upward in free space.'' That is
not a defect; it is general relativity showing up in a cheap chip.
\end{whybox}

\subsection{How MEMS makes it tiny}

A modern accelerometer is a \emph{MEMS} device (Micro-Electro-Mechanical System): the
mass and springs are etched out of silicon, tens of microns across. The proof mass
carries a comb of fingers interleaved with fixed fingers, forming a set of capacitors.
When the mass moves, one gap shrinks and the other grows, so the differential
capacitance changes by femtofarads. A capacitance-to-voltage converter turns this into
a voltage, an on-chip ADC digitizes it, and you read the number over I2C.

\begin{figure}[H]
\centering
\begin{tikzpicture}[scale=1.0,font=\small]
  \draw[thick] (-3.4,-1.4) rectangle (3.4,1.4);
  \draw[thick,decorate,decoration={zigzag,segment length=4pt,amplitude=3pt}]
     (-3.4,0) -- (-1.6,0);
  \draw[thick,decorate,decoration={zigzag,segment length=4pt,amplitude=3pt}]
     (1.6,0) -- (3.4,0);
  \draw[thick,fill=blue!12,rounded corners] (-1.6,-0.9) rectangle (1.6,0.9);
  \node at (0,0) {proof mass};
  \foreach \y in {-0.6,-0.2,0.2,0.6} {
     \draw[thick,blue!60!black] (1.6,\y) -- (2.3,\y);
  }
  \foreach \y in {-0.4,0,0.4} {
     \draw[thick,red!60!black] (3.4,\y) -- (2.5,\y);
  }
  \node[font=\footnotesize,align=center] at (2.3,1.85)
     {interleaved comb\\(differential capacitor)};
  \draw[-{Latex[]},thick] (-2.6,-1.85) -- (-1.0,-1.85)
     node[midway,below,font=\footnotesize] {frame accelerates $\to$};
  \draw[-{Latex[]},thick,blue!60!black] (0.2,-1.15) -- (-0.6,-1.15)
     node[midway,below,font=\footnotesize] {mass lags};
\end{tikzpicture}
\caption{A MEMS accelerometer in cartoon form. The proof mass hangs on silicon
springs; comb fingers form a differential capacitor whose imbalance reports the
mass's displacement, hence the specific force along this axis. Three such structures,
mutually perpendicular, give a 3-axis sensor.}
\label{ch34:fig-mems-accel}
\end{figure}

\subsection{Tilt sensing: getting angles from gravity}

Because a still accelerometer reports $-\vec g$ in body coordinates, it is a very good
inclinometer. If the chip is tilted, gravity's $9.81\,\mathrm{m/s^2}$ gets distributed
across the three axes, and the split tells you the orientation relative to vertical.
With $(a_x,a_y,a_z)$ the measured components (in any consistent unit), the standard
formulas are
\[
  \text{roll}\;\phi = \operatorname{atan2}(a_y,\,a_z), \qquad
  \text{pitch}\;\theta = \operatorname{atan2}\!\big(-a_x,\,\sqrt{a_y^2+a_z^2}\big).
\]

\begin{example}[Reading tilt from a still chip]
\label{ch34:ex-tilt}
Suppose a level board reads $(a_x,a_y,a_z)=(0,0,+1)\,g$. Now tip it forward by
$30^\circ$ about the $y$-axis. Gravity now projects as
$a_x = -\sin 30^\circ = -0.5\,g$ and $a_z = \cos 30^\circ = 0.866\,g$, so
$\theta = \operatorname{atan2}(0.5,\,0.866) = 30^\circ$. The accelerometer has become a
protractor -- \emph{as long as the chip is not also being shaken}, because any real
acceleration adds to gravity and corrupts the angle.
\end{example}

\begin{trapbox}
Tilt-from-gravity works only when the device is (nearly) still. The instant the arm
moves, $\vec a_{\text{true}}\ne 0$ and the measured vector is the \emph{sum} of gravity
and motion; you can no longer separate them from one sample. You also cannot get
\emph{yaw} (heading, rotation about vertical) from an accelerometer at all -- rotating
about the gravity vector leaves the gravity projection unchanged. Yaw needs the
gyroscope or the magnetometer.
\end{trapbox}

\section{The gyroscope: rotation through the Coriolis effect}
\label{ch34:gyro}

A MEMS gyroscope has no spinning wheel. Instead it exploits the \emph{Coriolis force}:
a mass driven to vibrate back and forth along one axis, when the whole frame rotates,
feels a sideways force proportional to the rotation rate. Concretely, a proof mass is
electrostatically driven to oscillate at velocity $\vec v$ along a ``drive'' axis. If
the chip rotates with angular rate $\vec\Omega$, the mass experiences a Coriolis
acceleration
\[
  \vec a_{\text{Cor}} \;=\; -\,2\,\vec\Omega \times \vec v,
\]
perpendicular to both. This tiny sideways motion is picked up by a second
(``sense'') capacitor, and its amplitude is proportional to $\abs{\vec\Omega}$
(Figure~\ref{ch34:fig-coriolis}). Three such resonators give angular rate about
$x$, $y$, $z$.

\begin{intuitionbox}
Picture walking straight out along a spoke of a spinning merry-go-round. Even though
you walk in a ``straight'' line, you get pushed sideways -- that push is Coriolis, and
its strength grows with how fast the platform spins. The MEMS gyro replaces your
walking with a mass vibrating thousands of times per second and measures the sideways
push electronically. No spin rate, no push; that is why a gyro reads zero when the chip
is held still, regardless of how it is oriented.
\end{intuitionbox}

\begin{figure}[H]
\centering
\begin{tikzpicture}[font=\small,scale=1.0]
  % drive oscillation (horizontal)
  \draw[thick,fill=blue!10,rounded corners] (-0.5,-0.5) rectangle (0.5,0.5);
  \node at (0,0) {\footnotesize mass};
  \draw[{Latex[]}-{Latex[]},very thick,blue!60!black] (-1.5,0) -- (-0.6,0);
  \draw[{Latex[]}-{Latex[]},very thick,blue!60!black] (0.6,0) -- (1.5,0);
  \node[blue!60!black,font=\footnotesize] at (0,-1.0) {drive $\vec v$ (vibration)};
  % rotation of frame
  \draw[-{Latex[]},very thick,structurepurple] (2.3,0.9) arc (60:-60:1.0);
  \node[structurepurple,font=\footnotesize] at (3.4,0) {$\vec\Omega$};
  % Coriolis (vertical) sense
  \draw[-{Latex[]},very thick,red!65!black] (0,0.6) -- (0,1.5);
  \draw[-{Latex[]},very thick,red!65!black] (0,-0.6) -- (0,-1.5);
  \node[red!65!black,font=\footnotesize,align=center] at (-1.9,1.2)
     {Coriolis push\\$-2\,\vec\Omega\times\vec v$\\(sense axis)};
\end{tikzpicture}
\caption{The Coriolis principle inside a MEMS gyroscope. The proof mass is driven to
vibrate along the horizontal drive axis (blue). When the frame rotates at rate
$\vec\Omega$ (purple), the mass feels a sideways Coriolis acceleration
$-2\,\vec\Omega\times\vec v$ (red) along the perpendicular sense axis; a sense capacitor
reports its amplitude, which is proportional to the rotation rate. No rotation means no
sideways push, so a still gyro reads zero.}
\label{ch34:fig-coriolis}
\end{figure}

\begin{focusbox}
\textbf{The gyro measures rate, not angle.} To get an \emph{orientation angle} you must
integrate: $\theta(t) = \theta(0) + \int_0^t \Omega(\tau)\,\dd\tau$. In discrete time,
$\theta[n] = \theta[n-1] + \Omega[n]\,\Delta t$. This integration is both the gyro's
superpower (smooth, fast, gravity-independent angle changes) and its curse (any
constant error in $\Omega$ accumulates without bound). We return to this as
\emph{drift} below.
\end{focusbox}

\begin{connectionbox}
Integrating angular rate to get angle is the discrete antiderivative from the calculus
companion: $\theta[n]=\theta[n-1]+\Omega[n]\Delta t$ is the rectangle rule applied
sample by sample. The same operation as a system has transfer function
$H(z) = \dfrac{\Delta t}{1-z^{-1}}$ -- a pure integrator, with a pole at $z=1$. From
Chapter~15 you already know a pole on the unit circle means infinite DC gain: the
integrator amplifies any constant (bias) input without limit. That single pole at
$z=1$ \emph{is} gyro drift, expressed in the language of the $z$-transform.
\end{connectionbox}

\section{The magnetometer: a chip-scale compass}
\label{ch34:mag}

A magnetometer measures the local magnetic flux density vector $\vec B$ in the chip's
frame. Most consumer parts use the \emph{Hall effect} or, more sensitively,
\emph{anisotropic/giant magnetoresistance} (AMR/GMR), where the electrical resistance
of a thin ferromagnetic film changes with the field threading it. Three orthogonal
sensing elements give $(B_x,B_y,B_z)$ in microtesla.

Outdoors and away from metal, the dominant field is the Earth's, magnitude roughly
$25$--$65\,\mathrm{\mu T}$ depending on latitude, tilted into the ground by the
\emph{inclination angle}. Projected onto the horizontal plane it points toward magnetic
north, which is why a magnetometer plus an accelerometer (to define horizontal) yields
a tilt-compensated compass -- the only one of the three sensors that gives an
\emph{absolute} heading rather than a relative change.

\begin{trapbox}
The magnetometer is the diva of the three. Any nearby ferromagnetic material or
current distorts the field: a \emph{hard-iron} offset (a permanent magnet or
magnetized screw shifts the whole reading by a constant vector) and a \emph{soft-iron}
distortion (nearby iron warps the field directionally, turning the locus of readings
from a sphere into an ellipsoid). A wrist device sits centimeters from a watch motor,
a metal buckle, maybe a phone -- magnetometer heading on the wrist is notoriously
unreliable and \emph{must} be calibrated (Chapter~35) before it is worth anything.
\end{trapbox}

\section{Units, ranges, and what the raw numbers mean}
\label{ch34:units}

Each sensor has a configurable \emph{full-scale range} (FSR). A wider range covers
more violent motion but spends fewer ADC codes per unit, so resolution drops. You read
a raw signed integer (typically 16-bit, $-32768$ to $+32767$) and convert with a
\emph{sensitivity} (LSB per physical unit) from the datasheet.

\begin{table}[H]
\centering
\renewcommand{\arraystretch}{1.2}
\begin{tabular}{@{}llll@{}}
\toprule
Sensor & Typical FSR options & Example sensitivity & Unit \\
\midrule
Accelerometer & $\pm2,\pm4,\pm8,\pm16\,g$ & $16384$ LSB/$g$ at $\pm2g$ & $g$ \\
Gyroscope & $\pm250,\pm500,\pm1000,\pm2000$ & $131$ LSB/(deg/s) at $\pm250$ & deg/s \\
Magnetometer & $\pm4800\,\mathrm{\mu T}$ (AK8963) & $0.15\,\mathrm{\mu T}$/LSB & $\mu$T \\
\bottomrule
\end{tabular}
\caption{Representative full-scale ranges and sensitivities for the MPU-9250 family.
The physical value is (raw integer)$\div$(sensitivity), or raw$\times$resolution for
the magnetometer.}
\label{ch34:tab-ranges}
\end{table}

\begin{example}[From raw counts to $g$]
\label{ch34:ex-counts}
At the $\pm2\,g$ setting the accelerometer gives $16384$ counts per $g$. A reading of
$a_z^{\text{raw}} = 16500$ converts to $16500/16384 = 1.007\,g$ -- a chip lying flat,
reading the expected $+1\,g$ of supported gravity, with a 0.7\% bias we would trim out
in calibration. A gyro reading of $\Omega_x^{\text{raw}} = 262$ at $\pm250$ deg/s
becomes $262/131 = 2.0\,\mathrm{deg/s}$, a slow drift you might mistake for real
rotation if you skipped bias removal.
\end{example}

\begin{focusbox}
\textbf{Choosing a range for the wrist.} Human limb motion rarely exceeds a few $g$ at
the wrist for ordinary gestures, but a sharp flick or a tap against a surface can spike
to $8$--$16\,g$. Wrist angular rates in fast gestures reach several hundred deg/s. For
a gesture controller, $\pm8\,g$ and $\pm1000\,\mathrm{deg/s}$ are sensible defaults:
wide enough not to clip on a snap, tight enough to keep resolution.
\end{focusbox}

\subsection{Noise, bias, and drift -- the three error families}

These three words get used loosely; keep them distinct, because each is fought with a
different tool.

\begin{itemize}
\item \textbf{Noise} is random, zero-mean, sample-to-sample jitter. It is quoted as a
  \emph{noise density}, e.g.\ $300\,\mathrm{\mu g/\sqrt{Hz}}$ for an accelerometer or
  $0.01\,\mathrm{(deg/s)/\sqrt{Hz}}$ for a gyro. Multiply by $\sqrt{\text{bandwidth}}$
  to get RMS noise. Noise is beaten by \emph{averaging / low-pass filtering}, at the
  cost of bandwidth (Chapter~20's SNR tradeoff exactly).
\item \textbf{Bias} (offset) is a constant added to every reading -- a gyro that reads
  $1.8\,\mathrm{deg/s}$ while perfectly still, or an accelerometer axis offset by
  $0.02\,g$. Bias is beaten by \emph{calibration}: measure it while stationary and
  subtract.
\item \textbf{Drift} is slow change \emph{over time} -- bias that wanders with
  temperature, or the unbounded angle error you get from integrating a biased gyro.
  Drift is the reason a gyro alone cannot hold orientation for more than seconds; it is
  beaten only by \emph{fusion} with an absolute reference (gravity, magnetic north) in
  Chapter~36.
\end{itemize}

\begin{connectionbox}
Gyro bias and noise interact viciously under integration. Treat the rate as
$\Omega[n] = \Omega_{\text{true}} + b + w[n]$ with constant bias $b$ and white noise
$w$ of per-sample variance $\sigma^2$. Integrating, the bias produces a \emph{linear}
angle error $b\,n\Delta t$, while the noise produces a random walk: the angle is the
sum $\theta_{\text{noise}}=\sum_{k=1}^n w[k]\,\Delta t$ of independent steps, so its
variance grows \emph{linearly} with the number of samples,
$\Var[\theta_{\text{noise}}] = n\,\Delta t^2\,\sigma^2$ (the discrete Brownian motion
of the probability companion). The standard deviation therefore grows like $\sqrt{n}$,
i.e.\ like $\sqrt{t}$. So the deterministic bias term ($\propto t$) dominates the
stochastic walk ($\propto \sqrt{t}$) at long times -- which is exactly why removing
bias buys you the most.
\end{connectionbox}

\begin{figure}[H]
\centering
\begin{tikzpicture}
\begin{axis}[
  width=12cm, height=5.4cm,
  xlabel={time $t$ (s)}, ylabel={integrated angle error (deg)},
  domain=0:10, samples=120,
  xmin=0, xmax=10, ymin=-2, ymax=22,
  axis lines=left, legend pos=north west, legend cell align=left,
  enlargelimits=false,
]
\addplot[specred, very thick] {1.8*x};
\addlegendentry{bias $1.8^\circ$/s $\rightarrow$ linear}
\addplot[sigblue, thick, dashed] {0.9*sqrt(x)};
\addlegendentry{noise $\rightarrow \sqrt{t}$ envelope}
\addplot[sigblue, thick, dashed, forget plot] {-0.9*sqrt(x)};
\end{axis}
\end{tikzpicture}
\caption{Why integrating a gyroscope drifts. A constant bias of $1.8^\circ$/s grows
the angle error \emph{linearly} (red), swamping the $\sqrt{t}$ random walk from white
noise (blue, dashed envelope). After ten seconds the bias alone has produced an
$18^\circ$ error. Remove bias first; fuse for the rest.}
\label{ch34:fig-drift}
\end{figure}

\section{Reading an IMU over I2C}
\label{ch34:reading}

The three sensors share a serial bus, almost always \emph{I2C} (two wires: SDA data,
SCL clock) or sometimes SPI for higher rates. You address the chip, point at a register,
and read back bytes. The mechanics are identical across the popular boards; only the
register map differs.

\begin{buildbox}
\textbf{Wiring an MPU-6050/9250 to an Arduino or Pi.} Four wires: \texttt{VCC}
(3.3\,V -- many breakouts have a regulator and tolerate 5\,V, but the bare chip is
3.3\,V), \texttt{GND}, \texttt{SDA}, \texttt{SCL}. On an Arduino Uno, \texttt{SDA}=A4,
\texttt{SCL}=A5; on a Raspberry Pi, GPIO2/GPIO3 (pins 3/5). Add $4.7\,\mathrm{k\Omega}$
pull-ups to 3.3\,V on SDA and SCL if your breakout lacks them. The default 7-bit
address is \texttt{0x68} (\texttt{0x69} if AD0 is tied high). \\[4pt]
On power-up the MPU-6050 is asleep; you must clear the \texttt{PWR\_MGMT\_1} register
(\texttt{0x6B}) to wake it. Accelerometer data starts at register \texttt{0x3B}, six
bytes (X,Y,Z high/low), big-endian; gyro at \texttt{0x43}. On the 9250 the AK8963
magnetometer hangs off an auxiliary I2C bus and needs the chip put in
``bypass/passthrough'' mode (or master mode) before you can reach it -- a classic
gotcha that sends beginners hunting for a dead magnetometer that was never enabled.
\end{buildbox}

\begin{lstlisting}[language=Python,caption={Minimal MPU-6050 read over I2C on a Raspberry Pi.},label=ch34:py-mpu]
import smbus, time
bus = smbus.SMBus(1)          # I2C bus 1 on a Pi
ADDR = 0x68
bus.write_byte_data(ADDR, 0x6B, 0)        # wake from sleep

def read_word(reg):
    hi = bus.read_byte_data(ADDR, reg)
    lo = bus.read_byte_data(ADDR, reg + 1)
    val = (hi << 8) | lo
    return val - 65536 if val >= 0x8000 else val   # signed 16-bit

ACC_LSB = 16384.0   # LSB per g at +/-2g
GYR_LSB = 131.0     # LSB per (deg/s) at +/-250 deg/s
while True:
    ax = read_word(0x3B) / ACC_LSB
    ay = read_word(0x3D) / ACC_LSB
    az = read_word(0x3F) / ACC_LSB
    gx = read_word(0x43) / GYR_LSB
    print(f"a=({ax:+.2f},{ay:+.2f},{az:+.2f}) g  gx={gx:+.1f} deg/s")
    time.sleep(0.02)   # ~50 Hz
\end{lstlisting}

\begin{buildbox}
\textbf{The BNO055 shortcut.} Bosch's BNO055 bundles an accelerometer, gyro,
magnetometer \emph{and} a built-in ARM core running sensor fusion. You can ask it
directly for an orientation quaternion or Euler angles over I2C -- no integration math
on your part. This is wonderful for a quick prototype and dangerous for understanding:
the fusion is a black box you cannot tune, and its magnetometer-driven yaw still
suffers all the wrist-proximity problems above. Prototype on a BNO055, but learn on a
raw MPU so you know what the black box is hiding (Chapter~36 opens it).
\end{buildbox}

\begin{focusbox}
\textbf{Sampling rate and aliasing -- the same Nyquist rule.} An IMU is an ADC, so
Chapter~18 applies without exception. If you read at $50\,\mathrm{Hz}$, any vibration
above $25\,\mathrm{Hz}$ aliases down into your gesture band and masquerades as slow
motion. Good IMUs include a configurable analog/digital low-pass (the
\texttt{CONFIG}/\texttt{DLPF} register on the MPU) -- set its cutoff below your output
Nyquist frequency. For wrist gestures, an internal DLPF around $20$--$40\,\mathrm{Hz}$
and an output rate of $100$--$200\,\mathrm{Hz}$ is a comfortable, alias-safe choice.
\end{focusbox}

\section{What wrist motion actually looks like}
\label{ch34:wrist-traces}

Enough abstraction -- what do the numbers \emph{do} when an arm moves? Three
qualitative regimes dominate any wrist trace, and learning to read them by eye is half
of gesture recognition (Chapter~38).

\begin{enumerate}
\item \textbf{Rest.} Accelerometer reads a steady vector of magnitude $\approx 1\,g$
  whose direction encodes wrist orientation; gyro reads $\approx 0$ plus bias and
  noise; magnetometer reads a steady (if distorted) field.
\item \textbf{Slow reorientation} (turning the wrist). Gyro shows a smooth bump
  (nonzero rate during the turn, back to zero after); the accelerometer's gravity
  split shifts to a new steady direction; integrating the gyro bump gives the angle
  turned.
\item \textbf{Sharp gesture} (a flick, tap, or snap). Both the accelerometer and gyro
  show a brief, large, oscillatory transient -- a spike with ringing as the arm
  decelerates. This is the signature gesture recognizers key on.
\end{enumerate}

\begin{figure}[H]
\centering
\begin{tikzpicture}
\begin{axis}[
  width=13cm, height=4.6cm,
  xlabel={time $t$ (s)}, ylabel={accel (g)},
  domain=0:6, samples=200,
  xmin=0, xmax=6, ymin=-2.2, ymax=2.6,
  axis lines=left, enlargelimits=false,
  title={Accelerometer (one axis) across a wrist sequence},
]
\addplot[sigblue, thick] {
  ( x<1.5 )*( 1 + 0.03*sin(deg(40*x)) )
+ ( x>=1.5 && x<3 )*( 1 - 0.55*(x-1.5)/1.5 )
+ ( x>=3 && x<3.8 )*( 0.45 + 0.03*sin(deg(40*x)) )
+ ( x>=3.8 && x<4.4 )*( 0.45 + 1.8*sin(deg(360*(x-3.8)/0.6))*exp(-3*(x-3.8)) )
+ ( x>=4.4 )*( 0.45 + 0.03*sin(deg(40*x)) )
};
\node[font=\footnotesize] at (axis cs:0.75,1.6) {rest $\approx 1g$};
\node[font=\footnotesize] at (axis cs:2.25,-1.4) {slow tilt};
\node[font=\footnotesize] at (axis cs:4.1,2.2) {flick};
\end{axis}
\end{tikzpicture}
\caption{A stylized single-axis accelerometer trace through a wrist sequence: a steady
$\approx 1\,g$ at rest, a smooth baseline shift as the wrist slowly tilts (gravity
re-projecting onto this axis), then a sharp damped-oscillatory transient from a flick.
Real traces are noisier but show these same three textures.}
\label{ch34:fig-wrist-accel}
\end{figure}

\begin{figure}[H]
\centering
\begin{tikzpicture}
\begin{axis}[
  width=13cm, height=4.6cm,
  xlabel={time $t$ (s)}, ylabel={rate (deg/s)},
  domain=0:6, samples=200,
  xmin=0, xmax=6, ymin=-260, ymax=260,
  axis lines=left, enlargelimits=false,
  title={Gyroscope (same axis, same sequence)},
]
\addplot[specred, thick] {
  ( x<1.5 )*( 2 + 1.5*sin(deg(55*x)) )
+ ( x>=1.5 && x<3 )*( 80*sin(deg(180*(x-1.5)/1.5)) )
+ ( x>=3 && x<3.8 )*( 2 + 1.5*sin(deg(55*x)) )
+ ( x>=3.8 && x<4.4 )*( 230*sin(deg(540*(x-3.8)/0.6))*exp(-3.5*(x-3.8)) )
+ ( x>=4.4 )*( 2 + 1.5*sin(deg(55*x)) )
};
\node[font=\footnotesize] at (axis cs:0.75,60) {rest $\approx 0$};
\node[font=\footnotesize] at (axis cs:2.25,150) {turn (one bump)};
\node[font=\footnotesize] at (axis cs:4.4,-180) {flick (ringing)};
\end{axis}
\end{tikzpicture}
\caption{The gyroscope on the same sequence. At rest it sits near zero (a small bias
visible). The slow turn is a single smooth bump whose \emph{area} is the angle turned.
The flick is a violent damped oscillation -- the rate signature recognizers exploit.}
\label{ch34:fig-wrist-gyro}
\end{figure}

\begin{happeningbox}
Notice the division of labor across the two plots. During the slow tilt the
accelerometer baseline \emph{shifts} (gravity re-projects) while the gyro shows a
\emph{bump} (the rate of that shift). During rest the accelerometer holds a nonzero
$1\,g$ value while the gyro sits at zero. The two sensors encode the \emph{same}
orientation change in complementary ways: the accelerometer sees it as a new static
direction, the gyro sees it as a transient rate. That complementarity is precisely what
makes fusion possible -- and it is why neither sensor alone suffices.
\end{happeningbox}

\section{What the wrist can and cannot read}
\label{ch34:limits}

This is the most important section in the chapter, and the one most often skipped in
hobbyist tutorials. We are building an instrument controller driven by \emph{finger
movements}, and we need a precise map of which channel carries which information. The
headline, stated up front so the rest of the chapter does not read as a letdown:
\emph{finger movement is detectable from a forearm/wrist band} -- just not from the IMU,
and not as continuous joint angles. It lives in the tendon and electrical channels,
which we read with the right sensor plus a per-finger calibration step.

\begin{goalbox}
\textbf{What the wrist band \emph{can} do for the instrument controller:}
\begin{itemize}
\item \textbf{IMU:} detect discrete gestures -- flicks, taps, twists, shakes, knocks,
  and the broad direction of an arm sweep -- which have large, distinctive
  accelerometer/gyro signatures
  (Figures~\ref{ch34:fig-wrist-accel}--\ref{ch34:fig-wrist-gyro}). Track wrist
  orientation (roll/pitch reliably from gravity; yaw only with fusion and a trustworthy
  magnetometer, which the wrist often is not), driving continuous parameters such as
  pitch-bend, volume, or a filter sweep. Track gross arm rhythm: conducting beats,
  strums, tempo from periodicity in the trace.
\item \textbf{FMG / EMG on the same band:} detect \emph{which finger moved and how
  hard}. The flexor and extensor tendons that drive the fingers run through the forearm,
  right under a wrist/forearm band. When you curl the index finger, a specific bundle of
  tendons bulges (a force-myography, FMG, signal you can read with a small pressure or
  strain array) and a specific group of motor units fires (a surface-EMG signal). Each
  finger produces a repeatable spatial \emph{signature} across the band
  (Figure~\ref{ch34:fig-tendon-channel}). With a short per-user calibration, those
  signatures become reliable per-finger detections.
\end{itemize}
\end{goalbox}

\begin{whybox}
Why fingers \emph{are} reachable from the forearm, even though the IMU is blind to them.
The muscles that move your fingers are not in your fingers -- they are in your forearm,
and they pull the fingers through long tendons. So the \emph{control signal} for every
finger physically passes under a forearm band twice over: mechanically, as tendon and
muscle bulk that an FMG pressure array feels, and electrically, as the motor-unit action
potentials that surface EMG picks up. The IMU misses all of this because finger muscles
move almost no mass at the wrist -- there is no inertial signature to see. Change the
sensor to one that reads tendons or electrical activity and the finger information is
right there, strong and repeatable. This is exactly what WristFlex (FMG) and the Myo /
MyoWare sEMG bands demonstrated.
\end{whybox}

\begin{figure}[H]
\centering
\begin{tikzpicture}[font=\small]
  % forearm cross-section
  \draw[thick,fill=orange!8] (0,0) circle (1.6);
  \node[font=\footnotesize] at (0,0) {forearm};
  % band ring of sensors
  \foreach \a in {0,45,...,315} {
     \fill[purple!55!black] (\a:1.6) circle (0.09);
  }
  \node[purple!55!black,font=\footnotesize] at (90:2.15) {FMG/EMG band};
  % tendon bundles inside, one highlighted
  \fill[blue!25] (-0.4,0.4) circle (0.18);
  \fill[blue!25] (0.4,0.5) circle (0.18);
  \fill[red!60!black] (0.6,-0.2) circle (0.2);
  \draw[-{Latex[]},red!60!black,thick] (2.6,-0.6) -- (0.85,-0.25)
     node[right=14pt,below,font=\footnotesize,align=left]
     {index-finger tendon\\bulges $\to$ local sensors fire};
  \node[blue!50!black,font=\footnotesize] at (-1.5,0.9) {other tendons};
\end{tikzpicture}
\caption{Why the fingers are readable from the forearm. A cross-section of the forearm
under the band: the muscles and tendons that drive each finger pass right beneath the
sensor ring. Moving one finger bulges a specific tendon bundle (red) and fires a
specific group of motor units, lighting up a characteristic subset of the FMG/EMG
sensors. Each finger leaves a repeatable spatial signature, which per-finger calibration
turns into a reliable detection. The IMU, by contrast, sits rigidly with the band and
feels none of this.}
\label{ch34:fig-tendon-channel}
\end{figure}

\begin{trapbox}
\textbf{The genuinely hard part, stated plainly:} a wrist-worn band cannot recover
\emph{continuous per-joint finger angles}. Telling \emph{that} the index finger flexed,
and roughly how forcefully, is easy and repeatable. Reconstructing the exact angle of
each of its three joints, continuously, is underdetermined: many distinct finger poses
produce nearly the same forearm tendon bulge and EMG pattern, so the mapping from
band signal to full hand pose is not invertible. The IMU is worse still -- it sees only
the bulk motion of the forearm and hand, so it carries \emph{no} per-finger information
at all. Continuous per-joint angles need sensors \emph{on the fingers}: IMU rings, flex
sensors, or instrumented gloves (see ``Learning on the Rings'' for finger-worn
continuous tracking). We give that as an upgrade path in the capstone, not as a
prerequisite -- because nothing the user wants for playing an instrument requires it.
\end{trapbox}

\begin{aestheticbox}
This framing is not a disappointment; it is the design brief. The IMU gives reliable
discrete gestures and wrist orientation; FMG/EMG on the same band gives reliable
per-finger detection after calibration; only continuous per-joint angles stay out of
reach without finger-worn sensors. That tells us exactly how to architect the
instrument: map wrist tilt and arm motion (IMU) to continuous expressive parameters,
read note and finger \emph{selection} from the muscle/tendon channel (FMG/EMG, with the
per-finger calibration of Chapter~37 and the capstone), and treat full continuous finger
pose as an optional ring-based upgrade. Every musical action the user actually wants --
which finger, how hard, with what wrist gesture and tilt -- sits on the doable side of
this line.
\end{aestheticbox}

\begin{historybox}
The path from spinning-wheel gyros to the chip on your wrist is short and recent.
Coriolis vibrating-structure gyros were proposed mid-20th century, but practical MEMS
versions only arrived in the 1990s as silicon micromachining matured (Charles Stark
Draper Laboratory and Bosch were early movers). The 2007 smartphone era drove the cost
of a full 6-axis IMU below a dollar, which is the only reason a hobbyist wrist tracker
is feasible at all. The Earth's magnetic field they lean on for heading has itself been
mapped for navigation since Gauss formalized geomagnetism in the 1830s -- the compass
in your IMU is the oldest idea in the package by two centuries.
\end{historybox}

\section{A worked end-to-end example}
\label{ch34:worked}

Let us tie units, integration, and drift together on one concrete number, the kind you
will actually compute in firmware.

\begin{example}[Estimating a wrist twist angle]
\label{ch34:ex-twist}
You sample the gyro's roll axis at $f_s = 100\,\mathrm{Hz}$ ($\Delta t = 0.01\,$s),
$\pm250\,\mathrm{deg/s}$ range ($131$ LSB per deg/s). During a one-second twist you
record raw values averaging $6550$ counts, i.e.\ $6550/131 = 50.0\,\mathrm{deg/s}$.
Integrating over the second:
\[
  \Delta\theta = \sum_{n=1}^{100} \Omega[n]\,\Delta t \approx 50.0 \times 1.0
  = 50^\circ .
\]
Now suppose you forgot to remove a bias of $1.8\,\mathrm{deg/s}$. Over the active
second it adds only $1.8^\circ$ -- tolerable. But if you keep integrating during the
following ten seconds of rest, the bias alone adds $18^\circ$ of phantom rotation
(Figure~\ref{ch34:fig-drift}). \emph{Lesson:} integrate the gyro only when motion is
actually happening, remove bias first, and lean on the accelerometer (gravity) to
re-anchor roll/pitch whenever the wrist is still -- the seed of complementary fusion.
\end{example}

\begin{labbox}
Wire up any IMU you have, log $(a_x,a_y,a_z,g_x,g_y,g_z)$ at $100\,\mathrm{Hz}$ to a
file, and perform four motions back to back: hold still 3\,s, slowly rotate the wrist
$90^\circ$, hold still 3\,s, then a single sharp flick. Plot all six channels. Confirm
by eye: (1) the accelerometer magnitude stays near $1\,g$ except during the flick;
(2) the slow rotation is a smooth gyro bump whose integral is about $90^\circ$;
(3) the flick is a brief, large transient in both sensors. You have now reproduced
Figures~\ref{ch34:fig-wrist-accel} and~\ref{ch34:fig-wrist-gyro} with your own arm.
\end{labbox}

\section{Summary}
\label{ch34:summary}

\begin{itemize}
\item An accelerometer measures \emph{specific force} $\vec a_{\text{true}}-\vec g$:
  it reads $+1\,g$ at rest, $0$ in free fall, and serves as a tilt sensor only when
  motion is small. It cannot sense yaw.
\item A gyroscope measures \emph{angular rate} via the Coriolis effect: zero at rest
  regardless of orientation, excellent for fast smooth angle \emph{changes}, but its
  integral drifts because constant bias accumulates linearly with time.
\item A magnetometer measures the \emph{magnetic field vector}, giving absolute heading
  -- but it is severely corrupted by nearby iron and currents, a real problem on the
  wrist, and demands calibration.
\item Raw readings are signed integers scaled by a datasheet sensitivity; choose a
  full-scale range wide enough not to clip and treat the IMU as an ADC bound by Nyquist
  (use the on-chip low-pass).
\item Errors come in three flavors -- noise (fight with filtering), bias (fight with
  calibration), drift (fight with fusion).
\item For the wrist controller: the IMU delivers discrete gestures and orientation, but
  \emph{not} individual fingers. Which finger moved, and how hard, is read instead from
  the tendon/electrical channel -- FMG or EMG on the same band, made reliable by a short
  per-finger calibration (Chapter~37 and the capstone). Only continuous per-joint finger
  angles stay out of reach without finger-worn sensors, an optional upgrade.
\end{itemize}

Chapter~35 takes the raw numbers seriously and builds proper calibration and sensor
models; Chapter~36 fuses the three sensors into a single drift-free orientation; and
Chapter~37 turns to the FMG/EMG sensors that read the fingers directly.

\section{Exercises}
\label{ch34:exercises}

\begin{exercise}
A 3-axis accelerometer at rest reads $(a_x,a_y,a_z)=(0.50,\,0.00,\,0.866)\,g$. What is
the chip's pitch and roll? Verify that the vector magnitude is consistent with the chip
being stationary, and explain what a magnitude of $1.3\,g$ would have told you instead.
\end{exercise}

\begin{exercise}
Explain, using Definition~\ref{ch34:def-specific-force}, why an accelerometer in a
satellite in orbit reads approximately zero on all axes even though the satellite is
very much under the influence of Earth's gravity.
\end{exercise}

\begin{exercise}
A gyroscope is set to $\pm2000\,\mathrm{deg/s}$, with sensitivity $16.4$ LSB per
deg/s. A raw reading is $-9840$ counts. Convert to deg/s. If this rate is held constant
for $0.5\,$s, what angle is swept?
\end{exercise}

\begin{exercise}
You integrate a gyro at $200\,\mathrm{Hz}$. The bias is $0.6\,\mathrm{deg/s}$ and you
fail to remove it. How much phantom angle accumulates after (a) 5\,s, (b) 1\,minute?
Compare to the $\sqrt{t}$ growth of a noise random walk with per-sample standard
deviation $0.05\,\mathrm{deg/s}$ and comment on which error dominates at one minute.
\end{exercise}

\begin{exercise}
At $\pm2\,g$ an accelerometer has $16384$ LSB/$g$; at $\pm16\,g$ it has $2048$
LSB/$g$. By what factor does the smallest resolvable acceleration change between the
two ranges, and what is the engineering tradeoff you are making when you widen the
range for the wrist controller?
\end{exercise}

\begin{exercise}
A magnetometer is rotated through all orientations while held far from metal; the cloud
of $(B_x,B_y,B_z)$ readings should trace a sphere. After mounting it on a wrist near a
steel buckle, the cloud becomes an off-center ellipsoid. Identify which distortion
shifts the center and which warps it into an ellipsoid, and state what each implies you
must estimate during calibration.
\end{exercise}

\begin{exercise}
You sample an IMU at $40\,\mathrm{Hz}$ with the internal low-pass disabled, and a
$55\,\mathrm{Hz}$ mechanical vibration (e.g.\ from a nearby motor) is present. To what
apparent frequency does it alias? Using Chapter~18's rule, choose an internal low-pass
cutoff and output rate that make a $30\,\mathrm{Hz}$ gesture band alias-free.
\end{exercise}

\begin{exercise}
Sketch (or describe) the accelerometer-magnitude and gyro traces you would expect from
each of: (a) the wrist held perfectly still; (b) a slow $180^\circ$ wrist rotation;
(c) two sharp taps a quarter-second apart. For each, name the feature a recognizer
would key on.
\end{exercise}

\begin{exercise}
A student claims their smartwatch app reads individual finger flexion ``from the IMU
alone.'' Using the physics of this chapter, give two concrete reasons the \emph{IMU}
cannot carry per-finger information, then name the sensor channel on the same band that
\emph{can} detect which finger moved, and explain in one sentence why the finger
control signal is physically present under a forearm band even though the IMU misses it.
\end{exercise}

\begin{exercise}
Detecting \emph{which} finger moved from a forearm band is achievable; reconstructing
the \emph{exact continuous angle} of every finger joint from the same band is not.
Explain the difference in terms of invertibility: why can many distinct finger poses
map to nearly the same FMG/EMG signature, and what kind of additional sensor removes the
ambiguity?
\end{exercise}

\begin{exercise}
Derive the discrete integrator transfer function $H(z)=\Delta t/(1-z^{-1})$ from the
update $\theta[n]=\theta[n-1]+\Omega[n]\Delta t$, and explain in one sentence how its
pole at $z=1$ corresponds to gyroscope drift.
\end{exercise}
